\subsubsection{Digit DP template}
\begin{lstlisting}[style=mystyle, language=C++]
// TC: O(D * 10 * state_size) where D is number of digits
// usa: contar numeros en [0, N] que cumplen una propiedad sobre sus digitos
#include <bits/stdc++.h>
using namespace std;

// skeleton only; adapt state, transition, and base case to your problem
string digitStr;
long long memo[20][2][200]; // dimensions depend on problem (pos, tight, state)

long long updateState(long long state, int d){
	// problem-specific; return new state after appending digit d
	return state;
}

long long digitDFS(int pos, bool tight, long long state){
	if(pos == (int)digitStr.size()) return 1; // adapt per problem (e.g., state matches)
	if(!tight && memo[pos][0][state] != -1) return memo[pos][0][state];
	int limit = tight ? digitStr[pos] - '0' : 9;
	long long res = 0;
	for(int d=0; d<=limit; d++){
		long long ns = updateState(state, d);
		res += digitDFS(pos + 1, tight && (d == limit), ns);
	}
	if(!tight) memo[pos][0][state] = res;
	return res;
}

long long digitDP(long long N){
	digitStr = to_string(N);
	memset(memo, -1, sizeof(memo));
	return digitDFS(0, true, 0);
}

int main(){
	// count numbers in [0, 999] whose digit sum equals 10
	digitStr = to_string(999);
	memset(memo, -1, sizeof(memo));
	cout << digitDFS(0, true, 0) << "\n";  // (skeleton; fill updateState)
	return 0;
}
\end{lstlisting}

\smallskip
\noindent\textbf{Versi\'on iterativa (forward).} Tabla \texttt{dp[pos][smaller][state]} llenada de \texttt{pos=0} a \texttt{pos=D}. \texttt{smaller=0} $\Rightarrow$ tight (prox d\'igito $\leq$ \texttt{N[pos]}); \texttt{smaller=1} $\Rightarrow$ libre (\texttt{0..9}).

\begin{lstlisting}[style=mystyle, language=C++]
// Cuenta cuantos x en [0, N) (= [0, N-1]) cumplen una propiedad sobre sus digitos.
// Para rango [a, b] INCLUSIVO:    digitDP(b+1) - digitDP(a).
// TC: O(D * 10 * (STATE_MAX + 1)).   Memoria: O(D * 2 * (STATE_MAX + 1)).

// Ejemplo de abajo: contar x sin dos digitos adyacentes iguales.
//   state = "ultimo digito colocado" (o NOT_STARTED al principio).
long long digitDP(long long N){
	// 1) digitos de N, mas significativo primero.
	vector<int> dig;
	if(N == 0){
		dig = {0};
	}else{
		long long t = N;
		while(t > 0){ dig.push_back((int)(t % 10)); t /= 10; }
		reverse(dig.begin(), dig.end());
	}
	int D = (int)dig.size();

	// 2) cota + sentinel. Para "no adyacentes iguales" alcanza 10.
	// adaptar, ahorita esta del 0 al 9 y el 10 de flag
	const int STATE_MAX   = 10;     
	const int NOT_STARTED = STATE_MAX;

	// 3) dp[pos][smaller][state], state en [0, STATE_MAX].
	vector<vector<vector<long long>>> dp(D + 1,
		vector<vector<long long>>(2, vector<long long>(STATE_MAX + 1, 0)));
	dp[0][0][NOT_STARTED] = 1;

	// 4) DP.
	for(int pos = 0; pos < D; pos++){
		for(int smaller = 0; smaller <= 1; smaller++){
			for(int digit = 0; digit <= STATE_MAX; digit++){
				long long ways = dp[pos][smaller][digit];
				// poda
				if(ways == 0) continue;                                 
				for(int nxt = 0; nxt <= (smaller ? 9 : dig[pos]); nxt++){
					// en este caso no puede haber 2 digitos consecutivos iguales
					// se puede modificar para se salten transiciones no validas
					if(nxt == digit) continue;
					
					int new_digit = nxt;
					if(digit == NOT_STARTED && nxt == 0){
						new_digit = NOT_STARTED;
					}
					dp[pos + 1][smaller || (nxt < dig[pos])][new_digit] += ways;
				}
			}
		}
	}

	// 5) Finales validos. dp[D][1][*] = "< N"; dp[D][0][*] = N exacto (no se suma).
	//    Para "no adyacentes": cualquier final vale (incluye NOT_STARTED -> el numero 0).
	long long ans = 0;
	for(int digit = 0; digit <= STATE_MAX; digit++)  
		ans += dp[D][1][digit];   // considera menores a N, [0,N)
	return ans;
}

int main(){
	// Ejemplo: contar x en [0, N) sin dos digitos adyacentes iguales.
	cout << digitDP(11)  << "\n";   // -> 11 (0..10 todos validos)
	cout << digitDP(12)  << "\n";   // -> 11 (sin el 11)
	cout << digitDP(100) << "\n";   // -> 91 
	return 0;
}
\end{lstlisting}
